\section{Problem Summary}
Rearrange the characters of the given uppercase string so that the result reads the same forward and backward. If no palindrome can be formed, print that clearly.

\section{Editorial}
A palindrome is controlled by character counts:
\begin{itemize}
\item every character used in the left half must also appear in the right half
\item therefore almost every count must be even
\item at most one character may have an odd count, because only the middle position can stay unmatched
\end{itemize}

So the algorithm is:
\begin{itemize}
\item count how many times each character appears
\item if more than one count is odd, the answer is impossible
\item otherwise build the left half from \(\lfloor \text{count} / 2 \rfloor\) copies of each letter
\item put the odd-count character, if it exists, in the middle
\item mirror the left half to obtain the right half
\end{itemize}

There is no need for search or backtracking. Once the counts are known, the structure of the palindrome is forced.

\section{Complexity Analysis}
\begin{itemize}
\item Time: \(O(n + \sigma)\), where \(\sigma\) is the alphabet size
\item Memory: \(O(\sigma)\)
\end{itemize}

\section{Notes}
For uppercase English letters, \(\sigma = 26\), so the counting array is tiny.
